\chapter{Materials and Methods}
\label{chap:methods}

% TODO: bridging paragraph outlining the experimental programme (Exp 0--6
%       from the project plan).

\section{Datasets}

\subsection{DS1 --- I-Support (20\,Hz)}
% TODO: 2-module / 6-chamber robot, 9 optical markers, 20\,Hz random-walk
%       pressure excitation, seeds 0--3, pickle layout (N x (27 + 6)).

\subsection{DS2 --- R01 recordings}
% TODO: single-segment, 3 valves (V1..V3), (X,Y,Z) in inches; 18{,}072
%       rows, likely time-series, units conversion.

\subsection{Data Splits and Preprocessing}
% TODO: DS1: split by seed (1,2 train / 3 val / 0x3 test); base-relative
%       tip coordinates (and mid for ablation); scaler persistence.
% TODO: DS2: temporal 70/15/15 split, NEVER shuffle.

\section{Exploratory Data Analysis (Exp 0)}
% TODO: workspace shape, marker-occlusion rate, time-series verification
%       for DS2, pressure--tip correlation.

\section{Analytical Baseline: PCC Inverse Kinematics}

\subsection{Single-Segment PCC (Exp 1A, DS2)}
% TODO: equations from Ma et al. (2022) eq. 11--21 (see
%       docs/references/pcc_equations.md).

\subsection{Two-Segment PCC (Exp 1B, DS1 --- stretch goal)}
% TODO: compound PCC + SQP optimisation (see
%       docs/references/multi_segment_pcc.md).

\section{Feed-Forward Neural IK}

\subsection{MLP Architecture (Exp 2A, Tip-Only)}
% TODO: inputs (X,Y,Z), hidden layers, activation, outputs = pressures.

\subsection{Tip + Mid Ablation (Exp 2B, DS1)}
% TODO: augment inputs with mid-segment marker mean; motivation: the
%       two-segment arm is underdetermined from tip alone.

\section{Temporal IK with LSTM (Exp 3)}
% TODO: windowing strategy, sequence length, stateful vs. stateless,
%       expected benefit on hysteresis.

\section{Jacobian-Learned IK (Exp 4)}
% TODO: learned-Jacobian + damped least-squares update rule from
%       Fang et al. (2022); see docs/references/jacobian_ik_algo.md.

\section{Simulation and Sim-to-Real with PyElastica (Exp 5)}

\subsection{Simulator Setup}
% TODO: Cosserat rod params from Manti 2016; pneumatic actuation model.

\subsection{Sim-to-Real Transfer Protocol}
% TODO: pre-train on synthetic, fine-tune on real DS2; evaluation protocol.

\section{Trajectory Tracking Evaluation (Exp 6)}
% TODO: unified protocol (figure-of-eight, circle, random) applied to
%       every IK method.

\section{Training Setup and Evaluation}

\subsection{Hardware, Software and Reproducibility}
% TODO: GPU, PyTorch version, seed = 42, checkpoint metadata layout.

\subsection{Hyperparameter Selection}
% TODO: validation-set driven search; never touch the test set.

\subsection{Metrics}
% TODO: MAE / RMSE / MAPE of tip position, trajectory error, inference
%       latency.
